% ─────────────────────────────────────────────────────────────────────────────
% Modèle de rapport d'activité — Évaluation à vague (CNRS, Comité national)
% Plan et consignes : « Recommandations pour la rédaction du dossier d'évaluation de l'activité des chercheurs et des chercheuses du CNRS », SGCN / DRH, novembre 2025.
% Compiler avec pdflatex (ou Overleaf). Complétez d'abord le bloc « À COMPLÉTER ».
% ─────────────────────────────────────────────────────────────────────────────
\documentclass[10pt,a4paper]{article}
\usepackage[utf8]{inputenc}
\usepackage[T1]{fontenc}
\usepackage[french]{babel}
\frenchsetup{StandardLists=true}
\usepackage[a4paper,top=2.3cm,bottom=2.2cm,left=2.1cm,right=2.1cm,headheight=14pt]{geometry}

% Police à empattement classique de LaTeX (Latin Modern), plus confortable à la lecture.
\usepackage{lmodern}
% Les consignes officielles demandent « Arial 10 » : pour s'y conformer strictement,
% remplacer la ligne ci-dessus par les deux suivantes.
% \usepackage[scaled=0.95]{helvet}
% \renewcommand{\familydefault}{\sfdefault}

\usepackage{microtype,xcolor,enumitem,titlesec,fancyhdr,lastpage,tabularx,comment}
\usepackage[most]{tcolorbox}
\usepackage{hyperref}

% ═══ À COMPLÉTER ═══════════════════════════════════════════════════════════
\newcommand{\Nom}{Prénom NOM}
\newcommand{\Grade}{CRCN}
\newcommand{\Unite}{UMR 0000 --- Nom du laboratoire}
\newcommand{\SectionCN}{Section 18 --- Chimie et vivant}
\newcommand{\Periode}{semestre 20.. --- semestre 20..}

% Consignes officielles (encadrés gris) : \consignestrue pour les afficher,
% \consignesfalse pour les masquer une fois le rapport rédigé.
\newif\ifconsignes
\consignestrue
% ═══════════════════════════════════════════════════════════════════════════

\definecolor{cnrs}{HTML}{22397F}
\definecolor{cnrsclair}{HTML}{EEF1FA}
\definecolor{gris}{HTML}{6B7280}
\hypersetup{colorlinks, linkcolor=cnrs, urlcolor=cnrs, pdftitle={Rapport d'activité}}

\setlength{\parindent}{0pt}
\setlength{\parskip}{0.45em}
\setlist{topsep=2pt,itemsep=1pt,leftmargin=1.3em}

% Intertitres : numéro en pastille, filet léger dessous
\titleformat{\section}{\large\bfseries\color{cnrs}}
  {\colorbox{cnrs}{\makebox[0.9em]{\color{white}\thesection}}}{0.6em}{}
  [{\color{cnrs!25}\titlerule[0.6pt]}]
\titlespacing*{\section}{0pt}{1.5em}{0.5em}
\titleformat{\subsection}{\normalsize\bfseries\color{cnrs}}{}{0em}{}
\titlespacing*{\subsection}{0pt}{1em}{0.3em}

% En-tête et pied de page
\pagestyle{fancy}
\fancyhf{}
\fancyhead[L]{\small\color{gris}Rapport d'activité --- \Nom}
\fancyhead[R]{\small\color{gris}Évaluation à vague}
\fancyfoot[C]{\small\color{gris}\thepage\,/\,\pageref{LastPage}}
\renewcommand{\headrulewidth}{0.3pt}
\renewcommand{\headrule}{{\color{gris!40}\hrule width\headwidth height\headrulewidth}}

% Consignes : encadré gris, ou rien du tout avec \consignesfalse
\ifconsignes
  \newtcolorbox{consigne}{enhanced, breakable, sharp corners, boxrule=0pt,
    colback=gris!7, borderline west={2pt}{0pt}{cnrs!45},
    left=8pt, right=6pt, top=4pt, bottom=4pt, before skip=4pt, after skip=8pt,
    fontupper=\small\itshape\color{gris!90!black}}
\else
  \excludecomment{consigne}
\fi

% Emplacement de votre texte
\newcommand{\arediger}{\textcolor{gris}{[À rédiger]}}

% Production scientifique : \production{référence}{lien HAL}{portée, impact, contribution}
\newcounter{production}
\newcommand{\production}[3]{%
  \refstepcounter{production}%
  \begin{tcolorbox}[enhanced, breakable, arc=2pt, boxrule=0.4pt, colframe=cnrs!25, colback=white,
      left=7pt, right=7pt, top=4pt, bottom=5pt, before skip=4pt, after skip=6pt]
    {\bfseries\color{cnrs}\theproduction.}~#1\par
    {\small\url{#2}}\par\smallskip
    #3
  \end{tcolorbox}}

\begin{document}

% ─── Titre ───────────────────────────────────────────────────────────────────
\begin{tcolorbox}[enhanced, sharp corners, boxrule=0pt, colback=cnrsclair,
    borderline west={4pt}{0pt}{cnrs}, left=14pt, right=10pt, top=10pt, bottom=10pt]
  {\footnotesize\color{cnrs}\MakeUppercase{Comité national de la recherche scientifique}}\par\vspace{3pt}
  {\LARGE\bfseries\color{cnrs}Rapport d'activité}\par\vspace{2pt}
  {\large Évaluation à vague --- 10 derniers semestres}\par\medskip
  {\small
  \begin{tabularx}{\linewidth}{@{}>{\bfseries}l X >{\bfseries}l X@{}}
    Nom & \Nom & Grade & \Grade \\
    Unité & \Unite & Section & \SectionCN \\
    Période & \Periode & & \\
  \end{tabularx}}
\end{tcolorbox}

\begin{consigne}
Période de référence : les 10 derniers semestres, indiquée dans le courriel d'ouverture de la campagne d'évaluation. Mise en forme : Arial 10, interligne 1. Conserver l'ordre des rubriques, même si certaines restent vides. Dépôt en PDF dans l'application e-valuation.
\end{consigne}


% ────────────────────────────────────────────────────────────────────────────
\section{Curriculum vitæ}
\begin{consigne}
Votre curriculum vitæ.
\end{consigne}
\arediger % Votre texte ici

% ────────────────────────────────────────────────────────────────────────────
\section{Résumé de votre activité}
\begin{consigne}
Résumé (1/2 page) qui permettra à l'ensemble des membres de la Section et/ou CID de prendre connaissance de votre profil scientifique.
\end{consigne}
\arediger % Votre texte ici

% ────────────────────────────────────────────────────────────────────────────
\section{Rapport d'activité}
\begin{consigne}
Vous présenterez de manière critique, en 30 pages maximum, les éléments suivants, en veillant à faire ressortir la valeur et l'influence de votre travail sur l'avancée des connaissances dans votre domaine :
\begin{itemize}
  \item votre activité de recherche au cours des 10 derniers semestres, les facteurs qui ont guidé vos choix, les résultats obtenus, leur portée et leur impact, et les éventuelles difficultés rencontrées ;
  \item votre participation à des collaborations avec des partenaires académiques, contractuelles ou non, dans un cadre français, européen ou international (préciser votre rôle, les partenaires, les thèmes des travaux et leur impact) ;
  \item vos participations à des contrats de recherche (ANR, contrats européens, autres agences) : intitulé long, acronyme, durée, partenaires et montants, nature de vos responsabilités (PI, co-PI, Work Package Manager), thème et contenu des travaux, portée et impact ;
  \item la place de votre recherche dans celle de votre unité, du point de vue des thématiques et des équipes (éventuellement sous forme d'organigramme) ;
  \item vos mobilités éventuelles, géographiques, thématiques ou fonctionnelles (séjours dans d'autres laboratoires publics ou privés, en France, en Europe ou à l'étranger, et leurs apports) ;
  \item les distinctions scientifiques reçues et les contributions qu'elles ont récompensées.
\end{itemize}
\end{consigne}
\arediger % Votre texte ici

\subsection*{Productions scientifiques (10 au maximum)}
\begin{consigne}
Classez vos productions par ordre d'importance. Pour chacune, expliquez sa portée, son impact et votre contribution personnelle, et donnez le lien vers la production elle-même dans HAL (ou une autre archive ouverte).
\end{consigne}
\production{Auteurs, « Titre », \emph{Revue}, volume (année), pages.}
  {https://hal.science/hal-00000000}
  {Portée, impact et contribution personnelle.}
\production{Auteurs, « Titre », \emph{Revue}, volume (année), pages.}
  {https://hal.science/hal-00000001}
  {Portée, impact et contribution personnelle.}

% ────────────────────────────────────────────────────────────────────────────
\section{Enseignement, formation et diffusion de la culture scientifique}
\begin{consigne}
Pour les 10 derniers semestres :
\begin{itemize}
  \item thèses dirigées et post-doctorats encadrés (noms, sujets, part prise dans l'encadrement, contenu et portée des travaux, devenir des docteurs et post-doctorants) ;
  \item participations à l'enseignement (niveau, titres des cours, établissements, nombre annuel d'heures) ;
  \item implication dans le partage, le management ou la publication des données de la recherche ;
  \item organisation de conférences, workshops, congrès et rencontres « portes ouvertes » ;
  \item revues ou ouvrages de vulgarisation et diffusion de l'information scientifique vers des publics non spécialistes ;
  \item interventions dans la presse écrite et audiovisuelle ;
  \item travaux d'expertise ;
  \item autres.
\end{itemize}
\end{consigne}
\arediger % Votre texte ici

% ────────────────────────────────────────────────────────────────────────────
\section{Transfert technologique, relations industrielles et valorisation}
\begin{consigne}
Pour les 10 derniers semestres :
\begin{itemize}
  \item contrats industriels et autres (durée, financeurs, partenaires et montants, nature de vos responsabilités, thème, contenu, portée et impact des travaux) ;
  \item projets de création d'entreprises ;
  \item travaux donnant lieu à des brevets ou à des développements de logiciels (votre rôle, contenu, difficultés rencontrées, impact) ;
  \item activités de consultance (destinataires, temps passé) ; distinctions reçues pour des activités de transfert ou de valorisation ;
  \item autres.
\end{itemize}
\end{consigne}
\arediger % Votre texte ici

% ────────────────────────────────────────────────────────────────────────────
\section{Encadrement, animation et management de la recherche}
\begin{consigne}
Pour les 10 derniers semestres, en mettant en lumière la manière d'exercer vos responsabilités, vos constats et les difficultés éventuelles rencontrées :
\begin{itemize}
  \item animation de programmes ou projets français, européens ou internationaux (type, objectif, ampleur, impact) ;
  \item direction d'équipe ou de laboratoire (nombre de personnes) ;
  \item autres responsabilités ou activités collectives (management de la recherche, missions d'intérêt général, conseils scientifiques, commissions, comités de lecture, etc.) ;
  \item autres.
\end{itemize}
\end{consigne}
\arediger % Votre texte ici

% ────────────────────────────────────────────────────────────────────────────
\section{Objectifs / Projet de recherche}
\begin{consigne}
Vos objectifs de recherche pour les 10 prochains semestres, notamment en perspective avec le projet de votre unité, et leur contribution aux grandes missions du CNRS.
\end{consigne}
\arediger % Votre texte ici

% ────────────────────────────────────────────────────────────────────────────
\section{Observations éventuelles}
\begin{consigne}
Observations que vous souhaitez porter à la connaissance de votre Section et/ou CID (1/2 page).
\end{consigne}
\arediger % Votre texte ici

\end{document}
